\section{Equations on products with a curve}\label{sec:curves}

The radial thresholds depend on dimension and geometry. A curve factor gives a useful comparison with exact one-variable primitives, also valid over a nonclosed Hermitian factor.

\begin{theorem}\label{thm:curve-products}\index{Curve!product equations}
Let $C$ be a compact Riemann surface with positive form $\omega_C$, let $(N,\omega_N)$ be a compact Hermitian $(n-1)$-fold, and put $h=\omega_C+\omega_N$ on $C\times N$. Let $p_0\in C$ and let $F\geq0$ be integrable, smooth positive away from $p_0$, with $\int_CF\omega_C=\int_C\omega_C$. There is a separable solution for the original measure $(F\circ\operatorname{pr}_C)h^n$ whose current has no divisor component on $\{p_0\}\times N$.

In a coordinate $z$ at $p_0$, write $r=|z|$ and $\omega_C=(\mathrm{i}/2\pi)g(z)\,\mathrm{d} z\wedge\mathrm{d}\bar z$.
\begin{enumerate}
\item If $Fg=C_0r^{2\tau}$, $C_0>0$ and $\tau>-1$, the total normal potential, modulo a smooth harmonic term, is $C_0r^{2(\tau+1)}/(\tau+1)^2$. It is bounded H\"older and minimal, while normal distance has order $r^{\tau+1}$. If $\tau=-a$ with $0<a<1$, its density is locally $L^p$ exactly when $pa<1$. Nonnegative integral vanishing orders give smooth potentials.
\item If $Fg=C_0r^{-2}t^{-b}$ with $t=\log(1/r)$, finite mass holds exactly for $b>1$. The atomless total normal potential is
\begin{equation}\label{eq:curve-log-primitive}
 P_b(t)=\begin{cases}
 \displaystyle\frac{4C_0t^{2-b}}{(b-1)(b-2)},&b\ne2,\\[4pt]
 -4C_0\log t,&b=2,
 \end{cases}
\end{equation}
modulo constants and smooth harmonic terms. It is bounded continuous and minimal precisely when $b>2$, and then is not H\"older. The normal end has finite distance precisely for $b>2$. For $1<b\leq2$, it is complete; the potential is unbounded and nonminimal but has zero Lelong number and trivial multiplier ideals. Every density in this logarithmic family lies outside $L^p$ for $p>1$.
\end{enumerate}
In the K\"ahler product case the unbounded solutions have full background mass, and closed uniqueness excludes a different bounded solution for the same measure. A small strict logarithmic barrier along $\{p_0\}\times N$ can coexist with the bounded minimal solutions.
\end{theorem}

\begin{proof}
On the curve, the real current $F\omega_C-\omega_C$ has integral zero and hence is $\ddc v$ for a global distribution $v$, unique up to a constant. Its quasi-subharmonic representative is smooth off $p_0$ by the local Poisson equation. Since $F\omega_C$ is atomless at $p_0$, its current has no atom. Pull $v$ back to $C\times N$. Positivity and the vanishing of any wedge with two curve factors give on the regular locus
\begin{equation}\label{eq:curve-product-equation}
 (h+\ddc v)^n=n(F\omega_C)\wedge\omega_N^{n-1}
 =(F\circ\operatorname{pr}_C)h^n.
\end{equation}
By \cref{thm:arbitrary-extended-ma}, the extended product agrees with this smooth product off $\{p_0\}\times N$ and gives that analytic set zero mass. The right side has the same property and finite mass. This proves the global equation, including when the curve potential has a finite discontinuous value at $p_0$.

For the power input put $P_\tau(r)=C_0r^{2(\tau+1)}/(\tau+1)^2$. Direct differentiation gives $\partial_z\bar\partial_zP_\tau=C_0r^{2\tau}$. Its radial boundary flux vanishes as $r\downarrow0$ when $\tau>-1$, so no Dirac term appears. The difference between the total potential and $P_\tau$ is harmonic on the full disk and smooth. Cut off the local potential and solve the remaining smooth mean-zero curve equation to realize the model globally. Its normal line element is comparable to $r^\tau|\mathrm{d} r|$, giving distance order $r^{\tau+1}$ by an all-path lower bound and a radial upper path. The $L^p$ assertion is the integral $\int_0^\epsilon r^{1-2pa}\,\mathrm{d} r$.

For the logarithmic input, $P_b''(t)=4C_0t^{-b}$ and $\partial_z\bar\partial_zP_b=P_b''/(4r^2)=C_0r^{-2}t^{-b}$. Its flux tends to zero for $b>1$, excluding a logarithmic Dirac mass. The same cutoff and smooth correction make a compact model. Density mass is $\int^\infty t^{-b}\,\mathrm{d} t$, finite exactly for $b>1$; its $p$th power for $p>1$ diverges after the logarithmic change of variables. The primitive has a finite limit exactly for $b>2$, with tail $t^{2-b}$, slower than every $r^\alpha$. Normal length is a constant times $\int^\infty t^{-b/2}\,\mathrm{d} t$, giving the same threshold. The projection to the curve bounds every path from below by that normal integral, while a fixed-$N$ radial path attains it.

For $1<b\leq2$, $|P_b(t)|/t\to0$, giving zero Lelong number. Every multiplier weight is bounded by an integrand of the form $\exp(-2t+o(t))$, which is integrable, so its multiplier ideals are trivial. These properties do not imply bounded ordinary type. In a K\"ahler product, the normalized curve integral gives full non-pluripolar background mass, and closed full-mass uniqueness excludes another bounded solution, as recalled in \cref{sec:history-regularity}. A bounded competitor for the original operator has deep truncations equal to itself globally, so the same exclusion applies there. Finally, for $D=\{p_0\}\times N$ and sufficiently small $\delta>0$, $h-\delta\Theta_D>0$; the potential $\delta\log|s_D|^2-C$ is a strict barrier.
\end{proof}
